% Generated draft: requires mathematical and source review.
\documentclass[11pt]{article}
\usepackage[margin=25mm]{geometry}
\usepackage{fontspec}
\setmainfont{Latin Modern Roman}
\usepackage{amsmath,amssymb,mathtools,unicode-math}
\setmathfont{Latin Modern Math}
\usepackage{enumitem,needspace,xurl,hyperref}
\setlength{\parindent}{0pt}
\setlength{\parskip}{5pt}
\setlength{\emergencystretch}{3em}
\newcommand{\N}{\mathbb{N}}
\newcommand{\Z}{\mathbb{Z}}
\newcommand{\Q}{\mathbb{Q}}
\newcommand{\R}{\mathbb{R}}
\newcommand{\C}{\mathbb{C}}
\newcommand{\F}{\mathbb{F}}
\newcommand{\A}{\mathbb{A}}
\newcommand{\PP}{\mathbb{P}}
\newcommand{\E}{\mathbb{E}}
\newcommand{\eps}{\varepsilon}
\newcommand{\dd}{\,\mathrm{d}}
\begin{document}
\section*{UnsolvedMath ID: 30001313}
\textbf{Source problem code:} catalogue entry
\subsection{Definitions and mathematical statement}
The supplied record poses an existence conjecture about a modification of Riemann-type minimal surfaces. Under the conventional interpretation that the ambient space is Euclidean \(\R^3\), an embedded surface \(M\subset\R^3\) is minimal if its mean curvature vanishes identically, and is properly embedded if its inclusion is an embedding whose inverse image of every compact set is compact. The ambient space is not explicitly specified in the supplied statement, so this interpretation requires confirmation. A planar end is understood informally as an end asymptotic to a plane; the precise asymptotic conditions are not supplied.

\textbf{Source-level conjecture.} There exists a one-parameter family of Riemann-like properly embedded minimal surfaces with one handle added between two consecutive planar ends.

Here ``Riemann-like'' is the original wording; the normalized statement uses ``Riemann-type.'' Neither term is defined in the supplied record. The intended family may be denoted \((M_t)_{t\in I}\), but the parameter domain \(I\), regularity of dependence on \(t\), and criterion for distinct family members are unspecified. The ordering that makes two planar ends consecutive and the precise topological meaning of adding one handle are also unspecified. In particular, the record does not determine whether the handle count concerns the entire surface, a fundamental region, or a quotient by a period. These qualifications remain unresolved; this draft preserves the source-level existence question without supplying missing hypotheses.
\subsection{Short English statement}
The conjecture asks for a one-parameter family of Riemann-like properly embedded minimal surfaces obtained by adding one handle between consecutive planar ends. The ambient setting, meaning of Riemann-like, family structure, and handle-counting convention require clarification before a fully precise formulation is possible.
\subsection{Sources}
\begin{itemize}
\item Supplied local UnsolvedMath registry record, permanent ID \(30001313\), title ``Riemann-Type Minimal Surfaces with One Handle,'' source problem code \texttt{catalogue entry}, in the Oberwolfach Reports Open Problems collection. The record gives proposed year 2009 and does not name a proposer.
\item Supplied report citation: \emph{Differentialgeometrie im Großen (2010)}, Oberwolfach Reports, pp.\ 1943--1972, DOI \url{https://doi.org/10.4171/owr/2009/35}. This is also the literature-source link labeled ``Differential geometry report.'' The citation's 2010 designation differs from the supplied proposed year 2009 and the year appearing in the DOI; their relationship requires checking.
\item Supplied literature-source link labeled ``Hauswirth--Pacard construction'': \url{https://doi.org/10.1007/s00222-007-0056-z}.
\item Supplied literature-source link labeled ``Morabito--Traizet construction'': \url{https://doi.org/10.1016/j.aim.2011.09.001}.
\item Literature assessment dated exactly \texttt{2026-08-21}: the supplied curation labels the problem ``partially solved'' and reports that higher-genus Riemann constructions produce closely related one-parameter families with handles, while the exact Riemann-like one-handle family in the extracted conjecture is not established in full. This is supplied dated curation derived from the OWR record, not an independently verified assessment; time-sensitive status and cited claims require independent checking.
\item The registry labels the statement \texttt{corrected\_verified} and reports that its mathematical content was checked against the official Oberwolfach report, with notation, context, and question form normalized without changing the source-level claim. This verification description is supplied metadata, not a verification performed for this draft.
\end{itemize}
\end{document}
